%!TEX TS-program = xelatex
%!TEX encoding = UTF-8 Unicode
%
\documentclass{gs-adonis}
\usepackage[english,italian]{babel}
%-------------------------------------------------------------------------------
%--------------------------------------------------------------------- CUSTOMS -
%-------------------------------------------------------------------------------
\usepackage[
  small,
  labelfont=bf,
  up,
  textfont=it,
  up
  ]{caption}
%
\usepackage{paralist}
\usepackage{subfigure}
\usepackage{csquotes}
\usepackage[style=ieee,backend=biber]{biblatex}
\bibliography{bibliografia.bib}
\usepackage{enumitem}
\usepackage{adjustbox}
\usepackage{hhline}
\usepackage{scrextend}
\usepackage{calc}
\usepackage{mwe}
\usepackage{todonotes}
\usepackage{hyperref}

% Font Alegreya (già incluso in gs-adonis, ma forziamo se necessario)
\usepackage[sfdefault]{AlegreyaSans}
\usepackage{Alegreya}
\renewcommand{\rmdefault}{Alegreya-LF}

%-------------------------------------------------------------------------------
%------------------------------------------------------------------------ MAIN -
%-------------------------------------------------------------------------------
\title{[INSERIRE TITOLO DEL PROGETTO]}
\subtitle{Progetto di \emph{Dottorato di Ricerca AFAM in Musica, Performance e Innovazione Tecnologica}}
\author{[NOME E COGNOME CANDIDATO/A]}
% secondary details
\correspondence{[indirizzo.email@esempio.com]}
\version{\today}
% headers
\runningauthor{[COGNOME CANDIDATO]}
\runningtitle{[TITOLO ABBREVIATO]}

%-------------------------------------------------------------------------------
%--------------------------------------------------------------------- COMANDI -
%-------------------------------------------------------------------------------
% Inserire qui eventuali comandi personalizzati
% Esempi:
% \newcommand{\opera}{\emph{Nome Opera}}
% \newcommand{\termine}{\emph{termine specifico}}

%-------------------------------------------------------------------------------
%-------------------------------------------------------------------- ABSTRACT -
%-------------------------------------------------------------------------------
\abstract{%
  % ABSTRACT DEL PROGETTO (massimo 300 parole)
  % Sintesi concisa del progetto di ricerca che evidenzi:
  % - Obiettivi principali
  % - Metodologie
  % - Risultati attesi
  % - Rilevanza per il campo di studi

  [Inserire qui l'abstract del progetto di ricerca, evidenziando gli aspetti chiave
  che lo collegano alle tematiche del dottorato: performance, artistic research,
  innovazione tecnologica, neuroscienze musicali, heritage veneziano, didattica
  speciale per l'inclusione.]
}

%-------------------------------------------------------------------------------
%-------------------------------------------------------------------- DOCUMENT -
%-------------------------------------------------------------------------------
\begin{document}
\maketitle

%-------------------------------------------------------------------------------
%------------------------------------------------------------------ KEYWORDS ---
%-------------------------------------------------------------------------------
\section*{Parole chiave}
% Massimo 5 parole chiave separate da virgole
[parola chiave 1], [parola chiave 2], [parola chiave 3], [parola chiave 4], [parola chiave 5].

%-------------------------------------------------------------------------------
%------------------------------------------------------ DESCRIZIONE PROGETTO ---
%-------------------------------------------------------------------------------
\section{Descrizione del Progetto di Ricerca}
{\small\textcolor{gray}{\textit{Limite totale: 10.000 caratteri spazi compresi}}}

%-------------------------------------------------------------------------------
%------------------------------------------------------ DESCRIZIONE SOGGETTO ---
%-------------------------------------------------------------------------------
\subsection{Descrizione del Soggetto di Ricerca}
{\small\textcolor{gray}{\textit{Caratteri suggeriti: 2.500--3.000 (25--30\% del totale)}}}

% In questa sezione descrivere:
% - Il campo di ricerca generale (performance studies, neuroscienze musicali, tecnologie digitali, etc.)
% - Il focus specifico del progetto
% - Le questioni teoriche e pratiche affrontate
% - L'inserimento nel panorama della ricerca musicale contemporanea
% - Collegamenti con l'heritage veneziano e le specificità culturali
% - Rilevanza per la didattica speciale e l'inclusione (se pertinente)

% Esempio di struttura:
% Il progetto si inserisce nell'ambito della ricerca su [area specifica],
% con particolare attenzione a [aspetto innovativo]. L'investigazione si
% concentra su [oggetto di studio] attraverso un approccio che integra
% [metodologie]. La ricerca si propone di esplorare [domanda principale]
% nel contesto specifico di [contesto veneziano/italiano/internazionale].

%-------------------------------------------------------------------------------
%--------------------------------------------------------- METODI E PROCESSO ---
%-------------------------------------------------------------------------------
\subsection{Metodi e Processo di Ricerca}
{\small\textcolor{gray}{\textit{Caratteri suggeriti: 3.000--3.500 (30--35\% del totale)}}}

% Descrivere dettagliatamente:
% - Metodologie di ricerca (artistic research, analisi musicale, etnomusicologia, etc.)
% - Strumenti e tecnologie utilizzate
% - Fonti primarie e secondarie
% - Eventuali collaborazioni con istituzioni, laboratori, centri di ricerca
% - Aspetti sperimentali e performativi
% - Tecniche di raccolta e analisi dati
% - Approcci interdisciplinari (neuroscienze, tecnologie, scienze sociali)

\subsubsection{Metodologie operative}
% Specificare le metodologie di ricerca artistica e scientifica

\subsubsection{Tecnologie e strumentazioni}
% Descrivere strumenti tecnologici, software, hardware necessari

\subsubsection{Fonti e materiali}
% Indicare archivi, repertori, documenti, registrazioni da consultare

\subsubsection{Collaborazioni previste}
% Eventuali partnership con enti, istituzioni, esperti

%-------------------------------------------------------------------------------
%------------------------------------------------------------ STATO DELL'ARTE ---
%-------------------------------------------------------------------------------
\subsection{Stato dell'Arte e Posizionamento della Ricerca}
{\small\textcolor{gray}{\textit{Caratteri suggeriti: 1.500--2.000 (15--20\% del totale)}}}

% Inquadrare il progetto rispetto alla letteratura esistente:
% - Studi precedenti nel campo specifico
% - Lacune nella conoscenza attuale
% - Originalità dell'approccio proposto
% - Dialogo con le ricerche internazionali
% - Riferimenti teorici principali

% Citare gli studi più rilevanti e spiegare come il progetto si posiziona
% rispetto al panorama della ricerca attuale, evidenziando gli elementi
% di novità e continuità.

%-------------------------------------------------------------------------------
%------------------------------------------------------- POSSIBILI RISULTATI ---
%-------------------------------------------------------------------------------
\subsection{Risultati Attesi e Impatto}
{\small\textcolor{gray}{\textit{Caratteri suggeriti: 2.000--2.500 (20--25\% del totale)}}}

% Descrivere i risultati previsti in termini di:

\subsubsection{Contributi alla conoscenza}
% Quale avanzamento si prevede rispetto allo stato attuale delle conoscenze

\subsubsection{Prodotti della ricerca}
% Specificare la tipologia di output:
% - Tesi di dottorato
% - Pubblicazioni scientifiche
% - Performance e concerti
% - Composizioni originali
% - Materiali didattici
% - Installazioni o opere multimediali
% - Altri prodotti artistici

\subsubsection{Distribuzione del lavoro triennale}
\textbf{Primo Anno:}
% Attività formative, ricerca preliminare, definizione metodologica

\textbf{Secondo Anno:}
% Sviluppo principale della ricerca, raccolta dati, sperimentazione

\textbf{Terzo Anno:}
% Completamento, sistematizzazione, produzione degli output finali

%-------------------------------------------------------------------------------
%-------------------------------------------------- RILEVANZA E DISSEMINAZIONE -
%-------------------------------------------------------------------------------
\subsection{Rilevanza per la Conoscenza, Comprensione e Pratica Musicale}
{\small\textcolor{gray}{\textit{Caratteri suggeriti: 1.000 (10\% del totale)}}}

% Spiegare l'importanza del progetto per:
% - La comunità scientifica e artistica
% - La didattica musicale
% - La pratica performativa
% - L'inclusione e la didattica speciale
% - Il patrimonio culturale veneziano e italiano
% - L'innovazione tecnologica in ambito musicale

\subsubsection{Disseminazione e condivisione}
% Modalità di condivisione dei risultati:
% - Conferenze e convegni
% - Pubblicazioni
% - Performance pubbliche
% - Workshop e masterclass
% - Piattaforme digitali
% - Collaborazioni internazionali

%-------------------------------------------------------------------------------
%-------------------------------------------------------------- BIBLIOGRAFIA ---
%-------------------------------------------------------------------------------
\section*{Bibliografia Iniziale}
{\small\textcolor{gray}{\textit{Non concorre al conteggio dei caratteri}}}

% Utilizzare lo stile bibliografico appropriato
% Includere le fonti principali per ogni area tematica:
% - Performance studies
% - Neuroscienze della musica
% - Tecnologie musicali
% - Heritage culturale
% - Didattica musicale e inclusione
% - Metodologie di ricerca artistica

\raggedright
\nocite{*}
\printbibliography

%-------------------------------------------------------------------------------
%---------------------------------------------------------------- APPENDICI ---
%-------------------------------------------------------------------------------
% \section*{Appendici}
% Eventuali materiali supplementari (grafici, spartiti, schemi tecnici, etc.)
% che non concorrono al conteggio dei caratteri

\end{document}
